\chapter{Sigma-Algebras and Measures}
\label{ch:iii01}

Measure theory begins by specifying which sets are stable under countable operations and which size functions are countably additive. The resulting continuity properties are the set-theoretic engine behind integration.

\section*{Learning goals}
After this chapter, the reader should be able to:
\begin{itemize}
\item verify that a proposed family of sets is a sigma-algebra and identify which closure axiom fails when it is not;
\item construct the sigma-algebra generated by elementary families and relate it to Borel sets in concrete spaces;
\item derive monotonicity and countable subadditivity from countable additivity;
\item apply continuity from below and continuity from above with the correct finiteness hypothesis;
\item distinguish finite, sigma-finite, complete, and incomplete measure spaces by examples;
\item construct null-set and completion examples showing why measurable structure matters beyond pointwise set inclusion;
\end{itemize}
\section*{Conceptual roadmap}
\[
\boxed{\text{structure}\;\longrightarrow\;\text{estimate}\;\longrightarrow\;\text{limit or representation}\;\longrightarrow\;\text{application}.}
\]

\paragraph{Mathematical setting.}
All sets to which $\mu$ is applied belong to a fixed sigma-algebra $\mathcal A$ on $X$; $\mu:\mathcal A\to[0,\infty]$. Continuity from above also holds if $\mu(E_N)<\infty$ at any stage, by applying the stated theorem to that tail. The completion consists of $A\cup N$ with $A\in\mathcal A$ and $N$ a subset of an $\mathcal A$-measurable null set, with measure $\mu(A)$.

\section{Sigma-algebras}
A sigma-algebra contains the empty set and is closed under complements and countable unions; De Morgan's laws give closure under countable intersections and differences.

\section{Generated sigma-algebras}
The sigma-algebra generated by a family is the intersection of all sigma-algebras containing it, giving the smallest measurable structure compatible with the generators.


% BEGIN VOL03-EXPANSION III01-example-01
\begin{example}[Generating a finite sigma-algebra]\label{ex:iii01-expand-01}
Let \(X=\{1,2,3,4\}\) and \(A=\{1,2\}\). Any sigma-algebra containing
\(A\) must contain \(A^c=\{3,4\}\), \(\varnothing\), and \(X\). These four
sets are already closed under complements and countable unions, so
\[
\sigma(A)=\{\varnothing,A,A^c,X\}.
\]
This proves both closure and minimality rather than merely guessing the list.
\end{example}
% END VOL03-EXPANSION III01-example-01
\section{Borel sets}
The Borel sigma-algebra is generated by open sets. On the real line, open intervals with rational endpoints already generate all Borel sets.

\section{Measures}
A measure is a nonnegative countably additive set function with value zero at the empty set. Countable additivity forces monotonicity and subadditivity.

\section{Finite and sigma-finite measures}
Finite measures have finite total mass; sigma-finite spaces are countable unions of finite-measure pieces and are the natural setting for product theory.

\section{Continuity from below}
If measurable sets increase to a union, their measures increase to the measure of that union.

\section{Continuity from above}
For decreasing measurable sets with finite measure at the first stage, measures decrease to the measure of the intersection.


% BEGIN VOL03-EXPANSION III01-example-02
\begin{example}[Why continuity from above needs finite mass]\label{ex:iii01-expand-02}
For Lebesgue measure on \(\mathbb R\), set \(E_n=[n,\infty)\). Then
\(E_n\downarrow\varnothing\), but \(m(E_n)=\infty\) for every \(n\).
Thus the measures do not decrease to \(m(\varnothing)=0\). The standard
continuity-from-above theorem needs \(\mu(E_1)<\infty\); this example isolates
the exact hypothesis used when the proof subtracts from \(\mu(E_1)\).
\end{example}
% END VOL03-EXPANSION III01-example-02
\section{Completion and null sets}
A complete measure space contains every subset of every measurable null set, so null-set modifications remain measurable.


% BEGIN VOL03-EXPANSION III01-example-03
\begin{example}[Completing a measure space]\label{ex:iii01-expand-03}
If \(N\) is measurable with \(\mu(N)=0\) and \(A\subset N\), the completion
declares \(A\) measurable with measure zero. More generally, completed sets
are obtained by adjoining subsets of measurable null sets to original
measurable sets. This is why changing a function on an arbitrary null subset
is safe only after completion of the measure space.
\end{example}
% END VOL03-EXPANSION III01-example-03
\section{Core structural results}

\begin{theorem}[Continuity from below]\label{thm:iii01-01}
If $E_n\uparrow E$, then $\mu(E_n)\uparrow\mu(E)$.
\begin{proof}
Disjointify the increase by writing $E$ as $E_1$ together with the successive differences $E_n\setminus E_{n-1}$, then use countable additivity.
\end{proof}
\end{theorem}

\begin{theorem}[Continuity from above]\label{thm:iii01-02}
If $E_n\downarrow E$ and $\mu(E_1)<\infty$, then $\mu(E_n)\downarrow\mu(E)$.
\begin{proof}
Apply continuity from below to $E_1\setminus E_n\uparrow E_1\setminus E$ and subtract from the finite quantity $\mu(E_1)$.
\end{proof}
\end{theorem}

\begin{theorem}[Countable subadditivity]\label{thm:iii01-03}
For measurable $E_n$, $\mu(\cup_nE_n)\le\sum_n\mu(E_n)$.
\begin{proof}
Replace the family by disjoint pieces $F_n=E_n\setminus\cup_{k<n}E_k$, use countable additivity, and compare $\mu(F_n)\le\mu(E_n)$.
\end{proof}
\end{theorem}

\section{Worked examples}

\begin{example}[The sigma-algebra generated by one subset]\label{ex:iii01-worked-01}
Let \(X=\{1,2,3,4\}\) and \(A=\{1,2\}\). Any sigma-algebra containing \(A\)
must contain \(A^c=\{3,4\}\), \(\varnothing\), and \(X\). These four sets are
closed under complements and countable unions, hence
\[
\sigma(A)=\{\varnothing,A,A^c,X\}.
\]
Forced membership gives minimality; direct closure shows no further sets are needed.
\end{example}

\begin{example}[Why continuity from above needs finite mass]\label{ex:iii01-worked-02}
For Lebesgue measure on \(\mathbb R\), let \(E_n=[n,\infty)\). Then
\[
E_n\downarrow\varnothing,\qquad m(E_n)=\infty\quad\text{for all }n.
\]
Thus \(m(E_n)\not\to m(\varnothing)\). The missing hypothesis is finite measure
at some stage of the decreasing sequence; without it the usual proof cannot
subtract from a finite reference mass.
\end{example}

\begin{example}[Completion adds subsets of null sets]\label{ex:iii01-worked-03}
Let \(N\in\mathcal A\) satisfy \(\mu(N)=0\), and let \(A\subset N\) be
nonmeasurable in the original sigma-algebra. In the completion, \(A\) is
declared measurable and has measure zero. More generally, completed measurable
sets have the form \(E\cup A\), with \(E\in\mathcal A\) and \(A\) contained in a
measurable null set.
\end{example}

\begin{example}[Countable subadditivity by disjointification]\label{ex:iii01-worked-04}
For measurable \(E_n\), set
\[
F_1=E_1,\qquad F_n=E_n\setminus\bigcup_{k<n}E_k.
\]
Then the \(F_n\) are disjoint, have the same union as the \(E_n\), and
\(F_n\subset E_n\). Hence
\[
\mu\!\left(\bigcup_nE_n\right)=\sum_n\mu(F_n)\le\sum_n\mu(E_n).
\]
\end{example}

\section{Solved dossiers}
Each dossier combines several ideas from the chapter in a longer argument. The aim is to make hypotheses, calculations, and proof strategy explicit.

\begin{problem}[The sigma-algebra generated by one subset]\label{prob:iii01-dossier-01}
Determine \(\sigma(A)\) for \(X=\{1,2,3,4\}\) and \(A=\{1,2\}\), proving both inclusion directions.
\end{problem}
\begin{solution}
\(\sigma(A)=\{\varnothing,A,A^c,X\}\). Every sigma-algebra containing \(A\)
must contain the four displayed sets, and those four already form a sigma-algebra.
\end{solution}

\begin{problem}[Why continuity from above needs finite mass]\label{prob:iii01-dossier-02}
Give a decreasing measurable sequence showing that continuity from above can fail when the initial mass is infinite.
\end{problem}
\begin{solution}
Take \(E_n=[n,\infty)\). Then \(E_n\downarrow\varnothing\), but
\(m(E_n)=\infty\) for every \(n\), whereas \(m(\varnothing)=0\).
\end{solution}

\begin{problem}[Completion adds subsets of null sets]\label{prob:iii01-dossier-03}
Explain concretely what new sets are added when a measure space is completed.
\end{problem}
\begin{solution}
Every subset of every measurable null set becomes measurable, with measure
zero; adjoining such subsets to old measurable sets gives the completed
sigma-algebra.
\end{solution}

\begin{problem}[Countable subadditivity by disjointification]\label{prob:iii01-dossier-04}
Derive countable subadditivity directly from countable additivity.
\end{problem}
\begin{solution}
Disjointify the family by \(F_n=E_n\setminus\cup_{k<n}E_k\), use countable
additivity on the \(F_n\), then use \(\mu(F_n)\le\mu(E_n)\).
\end{solution}

\begin{problem}[A finite sigma-algebra from a partition]\label{prob:iii01-dossier-05}
If \(X=A_1\sqcup A_2\sqcup A_3\), describe the sigma-algebra generated by the three atoms.
\end{problem}
\begin{solution}
It consists of all unions of the atoms \(A_1,A_2,A_3\), hence has \(2^3=8\)
sets (unless some atom is empty, in which case one first removes empty atoms).
\end{solution}

\begin{problem}[Sigma-finiteness of Lebesgue measure]\label{prob:iii01-dossier-06}
Show Lebesgue measure on \(\mathbb R\) is sigma-finite but not finite.
\end{problem}
\begin{solution}
\(\mathbb R=\bigcup_{n\ge1}[-n,n]\) and each interval has finite measure
\(2n\), while \(m(\mathbb R)=\infty\).
\end{solution}

\begin{problem}[Continuity from below]\label{prob:iii01-dossier-07}
Let \(E_n=(-n,n)\). Verify continuity from below for Lebesgue measure.
\end{problem}
\begin{solution}
The sets increase to \(\mathbb R\), and \(m(E_n)=2n\uparrow\infty=m(\mathbb R)\).
\end{solution}

\begin{problem}[Monotonicity from additivity]\label{prob:iii01-dossier-08}
Prove \(A\subset B\) implies \(\mu(A)\le\mu(B)\).
\end{problem}
\begin{solution}
Write \(B=A\sqcup(B\setminus A)\). Countable additivity gives
\(\mu(B)=\mu(A)+\mu(B\setminus A)\ge\mu(A)\).
\end{solution}

\begin{problem}[Continuity from below proof dossier]\label{prob:iii01-dossier-09}
Prove the following structural statement and identify the step where its main hypothesis is used: If $E_n\uparrow E$, then $\mu(E_n)\uparrow\mu(E)$.
\end{problem}
\begin{solution}
Disjointify the increase by writing $E$ as $E_1$ together with the successive differences $E_n\setminus E_{n-1}$, then use countable additivity. This proof also explains why the stated hypothesis is structural rather than cosmetic.
\end{solution}

\begin{problem}[Continuity from above proof dossier]\label{prob:iii01-dossier-10}
Prove the following structural statement and identify the step where its main hypothesis is used: If $E_n\downarrow E$ and $\mu(E_1)<\infty$, then $\mu(E_n)\downarrow\mu(E)$.
\end{problem}
\begin{solution}
Apply continuity from below to $E_1\setminus E_n\uparrow E_1\setminus E$ and subtract from the finite quantity $\mu(E_1)$. This proof also explains why the stated hypothesis is structural rather than cosmetic.
\end{solution}

\begin{problem}[Countable subadditivity proof dossier]\label{prob:iii01-dossier-11}
Prove the following structural statement and identify the step where its main hypothesis is used: For measurable $E_n$, $\mu(\cup_nE_n)\le\sum_n\mu(E_n)$.
\end{problem}
\begin{solution}
Replace the family by disjoint pieces $F_n=E_n\setminus\cup_{k<n}E_k$, use countable additivity, and compare $\mu(F_n)\le\mu(E_n)$. This proof also explains why the stated hypothesis is structural rather than cosmetic.
\end{solution}

\begin{problem}[Chapter synthesis]\label{prob:iii01-dossier-12}
Connect the main definitions and structural theorems of this chapter into one proof strategy. Explain what object is controlled, what estimate or closure principle is used, and what conclusion becomes available.
\end{problem}
\begin{solution}
Measure theory begins by specifying which sets are stable under countable operations and which size functions are countably additive. The resulting continuity properties are the set-theoretic engine behind integration. The recurring strategy is to encode the object in the chapter's natural measurable, normed, Fourier, or distributional structure, establish the controlling estimate, and then pass to the desired limit or representation.
\end{solution}
\section{Exercises with complete solutions}
The shorter exercise layer remains separate from the solved dossiers.

\begin{exercise}\label{exr:iii01-01}
State and apply the central principle behind Sigma-algebras. Give one immediate consequence in the setting of this chapter.
\end{exercise}
\begin{hint}
Check closure under complements and countable unions first; the remaining sigma-algebra operations follow from De Morgan's laws.
\end{hint}
\begin{solution}
A sigma-algebra contains the empty set and is closed under complements and countable unions; De Morgan's laws give closure under countable intersections and differences. As an immediate consequence, the same principle can be inserted into later convergence, approximation, transform, or weak-form arguments whenever its hypotheses are verified.
\end{solution}

\begin{exercise}\label{exr:iii01-02}
State and apply the central principle behind Generated sigma-algebras. Give one immediate consequence in the setting of this chapter.
\end{exercise}
\begin{hint}
Use the generated-sigma-algebra definition as an intersection of all sigma-algebras containing the generators.
\end{hint}
\begin{solution}
The sigma-algebra generated by a family is the intersection of all sigma-algebras containing it, giving the smallest measurable structure compatible with the generators. As an immediate consequence, the same principle can be inserted into later convergence, approximation, transform, or weak-form arguments whenever its hypotheses are verified.
\end{solution}

\begin{exercise}\label{exr:iii01-03}
State and apply the central principle behind Borel sets. Give one immediate consequence in the setting of this chapter.
\end{exercise}
\begin{hint}
For Borel measurability on \(\mathbb R\), reduce to intervals with rational endpoints rather than arbitrary open sets.
\end{hint}
\begin{solution}
The Borel sigma-algebra is generated by open sets. On the real line, open intervals with rational endpoints already generate all Borel sets. As an immediate consequence, the same principle can be inserted into later convergence, approximation, transform, or weak-form arguments whenever its hypotheses are verified.
\end{solution}

\begin{exercise}\label{exr:iii01-04}
State and apply the central principle behind Measures. Give one immediate consequence in the setting of this chapter.
\end{exercise}
\begin{hint}
Derive monotonicity by writing the larger set as a disjoint union of the smaller set and its difference.
\end{hint}
\begin{solution}
A measure is a nonnegative countably additive set function with value zero at the empty set. Countable additivity forces monotonicity and subadditivity. As an immediate consequence, the same principle can be inserted into later convergence, approximation, transform, or weak-form arguments whenever its hypotheses are verified.
\end{solution}

\begin{exercise}\label{exr:iii01-05}
State and apply the central principle behind Finite and sigma-finite measures. Give one immediate consequence in the setting of this chapter.
\end{exercise}
\begin{hint}
For sigma-finiteness, exhibit an explicit countable cover by measurable sets of finite measure.
\end{hint}
\begin{solution}
Finite measures have finite total mass; sigma-finite spaces are countable unions of finite-measure pieces and are the natural setting for product theory. As an immediate consequence, the same principle can be inserted into later convergence, approximation, transform, or weak-form arguments whenever its hypotheses are verified.
\end{solution}

\begin{exercise}\label{exr:iii01-06}
State and apply the central principle behind Continuity from below. Give one immediate consequence in the setting of this chapter.
\end{exercise}
\begin{hint}
Disjointify the increasing sequence and apply countable additivity term by term.
\end{hint}
\begin{solution}
If measurable sets increase to a union, their measures increase to the measure of that union. As an immediate consequence, the same principle can be inserted into later convergence, approximation, transform, or weak-form arguments whenever its hypotheses are verified.
\end{solution}

\begin{exercise}\label{exr:iii01-07}
State and apply the central principle behind Continuity from above. Give one immediate consequence in the setting of this chapter.
\end{exercise}
\begin{hint}
For continuity from above, subtract from the finite set \(E_1\); that is exactly where \(\mu(E_1)<\infty\) is used.
\end{hint}
\begin{solution}
For decreasing measurable sets with finite measure at the first stage, measures decrease to the measure of the intersection. As an immediate consequence, the same principle can be inserted into later convergence, approximation, transform, or weak-form arguments whenever its hypotheses are verified.
\end{solution}

\begin{exercise}\label{exr:iii01-08}
State and apply the central principle behind Completion and null sets. Give one immediate consequence in the setting of this chapter.
\end{exercise}
\begin{hint}
To check completeness, start with a measurable null set and ask whether every one of its subsets is measurable.
\end{hint}
\begin{solution}
A complete measure space contains every subset of every measurable null set, so null-set modifications remain measurable. As an immediate consequence, the same principle can be inserted into later convergence, approximation, transform, or weak-form arguments whenever its hypotheses are verified.
\end{solution}


% BEGIN VOL03-EXPANSION III01-exercises-01
\section{Graded supplementary exercises}
These exercises extend the preserved foundational set and are graded by mathematical role.

\subsection*{Standard computations}

\begin{exercise}[Generated family]\label{exr:iii01-expand-01}
On X={1,2,3}, find the sigma-algebra generated by {1}.
\end{exercise}
\begin{hint}
Close under complements and unions.
\end{hint}
\begin{solution}
It is {empty, {1}, {2,3}, X}.
\end{solution}

\begin{exercise}[Continuity below]\label{exr:iii01-expand-02}
If E\_n increases to E and mu(E\_n)=1-2\textasciicircum{}(-n), find mu(E).
\end{exercise}
\begin{hint}
Take the monotone limit.
\end{hint}
\begin{solution}
mu(E)=1.
\end{solution}

\begin{exercise}[Subadditivity]\label{exr:iii01-expand-03}
If mu(E)=2 and mu(F)=4, give a universal upper bound for mu(E union F).
\end{exercise}
\begin{hint}
Use subadditivity.
\end{hint}
\begin{solution}
The bound is 6.
\end{solution}

\begin{exercise}[Sigma-finite line]\label{exr:iii01-expand-04}
Explain why Lebesgue measure on the real line is sigma-finite.
\end{exercise}
\begin{hint}
Cover by bounded intervals.
\end{hint}
\begin{solution}
The real line is the union of [-n,n], each having finite measure.
\end{solution}

\begin{exercise}[Null completion]\label{exr:iii01-expand-05}
If A is a subset of a measurable null set in a complete space, what is mu(A)?
\end{exercise}
\begin{hint}
Use completeness and monotonicity.
\end{hint}
\begin{solution}
A is measurable and has measure 0.
\end{solution}

\subsection*{Proofs}

\begin{exercise}[Monotonicity]\label{exr:iii01-expand-06}
Prove A subset B implies mu(A)<=mu(B).
\end{exercise}
\begin{hint}
Write B as a disjoint union.
\end{hint}
\begin{solution}
B=A disjoint union (B minus A), so nonnegativity gives the inequality.
\end{solution}

\begin{exercise}[Continuity below]\label{exr:iii01-expand-07}
Prove continuity from below.
\end{exercise}
\begin{hint}
Disjointify the increasing sequence.
\end{hint}
\begin{solution}
Write the union as E1 plus successive differences and use countable additivity.
\end{solution}

\begin{exercise}[Continuity above]\label{exr:iii01-expand-08}
Derive continuity from above from continuity below under finite first measure.
\end{exercise}
\begin{hint}
Apply continuity below to E1 minus En.
\end{hint}
\begin{solution}
Subtract the increasing limit from the finite number mu(E1).
\end{solution}

\begin{exercise}[Subadditivity proof]\label{exr:iii01-expand-09}
Prove countable subadditivity.
\end{exercise}
\begin{hint}
Replace sets by disjoint new pieces.
\end{hint}
\begin{solution}
Disjointify and compare each new piece with the original set.
\end{solution}

\subsection*{Counterexamples and hypothesis testing}

\begin{exercise}[Finite-mass necessity]\label{exr:iii01-expand-10}
Give a decreasing sequence showing continuity from above can fail without finite first measure.
\end{exercise}
\begin{hint}
Use tails of the real line.
\end{hint}
\begin{solution}
[n,infinity) decreases to empty but every term has infinite measure.
\end{solution}

\begin{exercise}[Algebra not sigma-algebra]\label{exr:iii01-expand-11}
Give a family closed under finite operations but not countable unions.
\end{exercise}
\begin{hint}
Use finite/cofinite subsets of the natural numbers.
\end{hint}
\begin{solution}
The union of even singleton sets is neither finite nor cofinite.
\end{solution}

\begin{exercise}[Incomplete measure]\label{exr:iii01-expand-12}
Describe what failure of completeness means.
\end{exercise}
\begin{hint}
Look inside measurable null sets.
\end{hint}
\begin{solution}
A subset of a measurable null set can fail to be measurable.
\end{solution}

\subsection*{Applications and investigations}

\begin{exercise}[Probability events]\label{exr:iii01-expand-13}
Why does countable additivity matter for mutually exclusive probabilistic outcomes?
\end{exercise}
\begin{hint}
Union their events.
\end{hint}
\begin{solution}
It makes probability of a countable disjoint union equal the sum of probabilities.
\end{solution}

\begin{exercise}[Null modifications]\label{exr:iii01-expand-14}
Why does completion matter for almost-everywhere arguments?
\end{exercise}
\begin{hint}
Exceptional sets can be arbitrary subsets of null sets.
\end{hint}
\begin{solution}
Completion keeps all such exceptional subsets measurable and null.
\end{solution}

\subsection*{Challenge problems}

\begin{exercise}[Rational generators]\label{exr:iii01-expand-15}
Explain why rational-endpoint intervals generate all Borel subsets of the real line.
\end{exercise}
\begin{hint}
Approximate open sets by countably many rational intervals.
\end{hint}
\begin{solution}
Every open set is a countable union of such intervals, so all Borel sets are generated.
\end{solution}

\begin{exercise}[Borel versus Lebesgue]\label{exr:iii01-expand-16}
Why is the completed Lebesgue sigma-algebra larger than the Borel sigma-algebra?
\end{exercise}
\begin{hint}
Add subsets of Borel null sets.
\end{hint}
\begin{solution}
The Cantor set is Borel and null and has cardinality $\mathfrak c$. It has $2^{\mathfrak c}$ subsets, whereas there are only $\mathfrak c$ Borel sets in $\mathbb R$ (each has a countable Borel code). Thus some of its subsets are non-Borel, and all belong to the completion. The Borel-code cardinality fact is a standard set-theoretic fact used here.
\end{solution}
% END VOL03-EXPANSION III01-exercises-01
\section*{Chapter summary}
The chapter now supplies a canonical theorem-and-dossier layer for subsequent measure, Fourier, distributional, Sobolev, and PDE arguments.
